\section{Results}
\label{sec:results}

\subsection{Decode speed at short context}

Figure~\ref{fig:speed} ranks the configurations on the six short prompts. Raising the
baseline's MTP drafts from two to five gains 19\%; beyond five, rejected drafts cost more than
they bring. OpenJet, with four MTP drafts on the baseline's own GGUF, reaches 134.7\,tok/s
(1.50$\times$). NInfer with MTP is slightly faster (137--140\,tok/s) and with seven DFlash2
drafts reaches 166.9\,tok/s (1.86$\times$). azllm, with seven DFlash2 drafts on its smaller
\texttt{TQ\_L} weights, reaches 228.6\,tok/s: 2.55$\times$ the baseline and 1.37$\times$
NInfer.

\begin{figure}[t]
  \centering
  \begin{tikzpicture}
    \begin{axis}[
        xbar, width=0.68\columnwidth, height=6.9cm, bar width=7pt,
        xmin=0, xmax=265, xlabel={decode speed (tok/s)},
        symbolic y coords={b2,b3,b4,b5,b6,oj,nm3,nm4,nm3e,nd7,az},
        ytick={b2,b3,b4,b5,b6,oj,nm3,nm4,nm3e,nd7,az}, y dir=reverse,
        yticklabels={llama.cpp MTP 2 (baseline), llama.cpp MTP 3, llama.cpp MTP 4,
                     llama.cpp MTP 5, llama.cpp MTP 6, OpenJet MTP 4, NInfer MTP 3 int8,
                     NInfer MTP 4 int8, NInfer MTP 3 E8, NInfer DFlash2 7 E8, azllm DFlash2 7 q8/q4},
        yticklabel style={font=\scriptsize}, xticklabel style={font=\scriptsize},
        xlabel style={font=\small},
        nodes near coords, nodes near coords style={font=\scriptsize},
        every axis plot/.append style={draw=none, bar shift=0pt},
        xmajorgrids, grid style={gray!20}]
      \addplot[fill=cbase] coordinates {(89.6,b2) (98.3,b3) (105.8,b4) (106.6,b5) (104.7,b6)};
      \addplot[fill=cjet] coordinates {(134.7,oj)};
      \addplot[fill=cnin2] coordinates {(137.2,nm3) (139.5,nm4) (140.5,nm3e)};
      \addplot[fill=cnin] coordinates {(166.9,nd7)};
      \addplot[fill=cgood] coordinates {(228.6,az)};
    \end{axis}
  \end{tikzpicture}

  \caption{Mean decode speed over six short prompts, recommended sampling, thinking off.
           Grey: llama.cpp; orange: OpenJet; blue: NInfer; green: azllm.}
  \label{fig:speed}
\end{figure}

The gain depends on the text (Figure~\ref{fig:perprompt}). On the two coding prompts and the
word problem, azllm runs 2.9--3.1$\times$ faster than the baseline and NInfer 2.1--2.3$\times$;
on free prose, which no drafter predicts well, the gains fall to 1.70$\times$ and
1.28$\times$. Thinking mode (effort \emph{medium}) gives 208.4 against 84.2\,tok/s for azllm
(2.47$\times$) and 153.3 for NInfer (1.82$\times$).

\begin{figure}[t]
  \centering
  \begin{tikzpicture}
    \begin{axis}[
        ybar, width=\columnwidth, height=5.2cm, bar width=4pt,
        ymin=0, ymax=320, ylabel={tok/s}, ylabel style={font=\small},
        symbolic x coords={code,code2,math,list,explain,prose},
        xtick=data, xticklabel style={font=\scriptsize}, yticklabel style={font=\scriptsize},
        legend style={font=\scriptsize, at={(0.5,1.03)}, anchor=south, draw=none, legend columns=2, /tikz/every even column/.append style={column sep=6pt}},
        legend cell align=left, enlarge x limits=0.1,
        every axis plot/.append style={draw=none}, ymajorgrids, grid style={gray!20}]
      \addplot[fill=cbase] coordinates {(code,97.0) (code2,94.7) (math,96.6) (list,94.6) (explain,84.0) (prose,70.5)};
      \addplot[fill=cjet] coordinates {(code,165.4) (code2,160.5) (math,150.7) (list,120.6) (explain,121.1) (prose,90.1)};
      \addplot[fill=cnin] coordinates {(code,212.9) (code2,214.7) (math,205.7) (list,151.2) (explain,132.3) (prose,90.1)};
      \addplot[fill=cgood] coordinates {(code,302.7) (code2,282.6) (math,276.1) (list,207.7) (explain,193.1) (prose,120.0)};
      \legend{Baseline, OpenJet MTP 4, NInfer DFlash2 7, azllm DFlash2 7}
    \end{axis}
  \end{tikzpicture}

  \caption{Decode speed per prompt. Structured output (code, a worked problem, a table) gains
           most; a short story gains least.}
  \label{fig:perprompt}
\end{figure}

\paragraph{Number of drafts.}
Figure~\ref{fig:drafts} sweeps the number of draft tokens. NInfer's DFlash2 peaks at five to
seven drafts and drops sharply beyond; azllm's DFlash2 rises up to seven and stays flat up to
ten. MTP, which drafts one token per step, peaks at three to five. On NInfer the KV format makes
no measurable difference at short context (DFlash2 7: 166.9\,tok/s with E8, 166.7 with int8);
on azllm neither (228.5 with \texttt{q8/q4}, 230.7 with \texttt{q4/q4}).

\begin{figure}[t]
  \centering
  \begin{tikzpicture}
    \begin{axis}[
        width=0.9\columnwidth, height=5.2cm, xmin=1.5, xmax=10.5, ymin=80, ymax=245,
        xlabel={draft tokens per step}, ylabel={tok/s},
        xlabel style={font=\small}, ylabel style={font=\small},
        ticklabel style={font=\scriptsize}, xtick={2,3,...,10},
        legend style={font=\scriptsize, at={(0.5,1.03)}, anchor=south, draw=none, legend columns=3, /tikz/every even column/.append style={column sep=4pt}},
        legend cell align=left, grid=major, grid style={gray!20}]
      \addplot[cbase, thick, mark=*] coordinates {(2,89.6) (3,98.3) (4,105.8) (5,106.6) (6,104.7)};
      \addplot[cnin2, thick, mark=square*] coordinates {(3,137.2) (4,139.5) (5,132.5)};
      \addplot[cnin, thick, mark=triangle*] coordinates {(3,136.3) (5,160.3) (6,155.6) (7,162.8) (8,136.8) (10,133.6)};
      \addplot[cjet, only marks, mark=diamond*, mark size=3pt] coordinates {(4,134.7)};
      \addplot[cgood, thick, mark=pentagon*] coordinates {(5,205.6) (6,213.1) (7,230.4) (8,230.1) (10,230.7)};
      \legend{llama.cpp MTP, NInfer MTP, NInfer DFlash2, OpenJet MTP, azllm DFlash2}
    \end{axis}
  \end{tikzpicture}

  \caption{Decode speed against the number of draft tokens. NInfer DFlash2 points use E8 KV,
           except 3 and 10 drafts (int8 KV at 96K); azllm points use \texttt{q8/q4} KV at 128K
           without vision.}
  \label{fig:drafts}
\end{figure}

\subsection{Accuracy}

Table~\ref{tab:profiles} lists accuracy next to speed for the retained configurations. With
greedy decoding and thinking off, all engines solve 238--244 of the 250 GSM8K problems and
156--158 of the 164 HumanEval problems. OpenJet, which reads exactly the baseline's weights,
solves four GSM8K problems fewer than the baseline, and azllm solves more with its 4-bit cache
(242) than with its 8-bit one (238), so differences of this size are the noise of greedy decoding
on near-ties. To compare the weights themselves we turn to perplexity.

\PROFILETABLE

\subsection{Weight precision}

Figure~\ref{fig:weights} compares the three weight files on the same 98{,}304 positions of code.
Taking the baseline's UD-Q4\_K\_XL as the reference (0.4330 nats per token), azllm's
\texttt{TQ\_L} loses 0.46\% and NInfer's groupwise artifact 1.78\%. The ranking holds in every
chunk. \texttt{TQ\_L} is the smallest file (4.31 against about 5.2 bits per weight), so its
rotated format buys both speed and precision. Prefilling with FP16 instead of int8 activations
changes azllm's result by 0.0002 nats, so the gap comes from the weights, not the prefill path.

\begin{figure}[t]
  \centering
  \begin{tikzpicture}
  \begin{axis}[
      xbar, width=0.58\columnwidth, height=3.6cm, bar width=9pt,
      xmin=0, xmax=2.6, xlabel={$\Delta$ NLL per token (\%)},
      xlabel style={font=\small},
      symbolic y coords={ud,tq,ni}, ytick={ud,tq,ni}, y dir=reverse,
      yticklabels={baseline, azllm, NInfer},
      yticklabel style={font=\scriptsize}, xticklabel style={font=\scriptsize},
      point meta=explicit symbolic, nodes near coords, nodes near coords style={font=\scriptsize, anchor=west},
      every axis plot/.append style={draw=none, bar shift=0pt}, xmajorgrids, grid style={gray!20}]
    \addplot[fill=cbase] coordinates {(0,ud) [0.4330 nats]};
    \addplot[fill=cgood] coordinates {(0.46,tq) [+0.46\%]};
    \addplot[fill=cnin] coordinates {(1.78,ni) [+1.78\%]};
  \end{axis}
\end{tikzpicture}

  \caption{Mean NLL per token on 98{,}304 positions of C++/CUDA code (32K--64K tokens of
           context each), relative to the baseline's weights; 8-bit KV in every engine. Lower is
           better.}
  \label{fig:weights}
\end{figure}

\subsection{Context capacity}

Figure~\ref{fig:capacity} maps the largest context each configuration accepts. Three factors
decide it: the KV format (int8 to rk8v4 to E8 roughly doubles NInfer's capacity), the drafter
(NInfer's DFlash2 keeps its own KV cache and costs 57--111K tokens), and the vision tower
(16--33K tokens). OpenJet allocates a fixed 2048-token prefill scratch; shrinking it to 512
tokens raises its ceiling from 131{,}072 to 147{,}456 tokens with \texttt{q8} keys and values,
and it has no 4-bit KV and no vision. azllm reaches the full native window, with vision, with
\texttt{q4/q4} KV, and 241{,}664 tokens with \texttt{q8/q4}. Near its limits it moves the token
embeddings and a copy of the LM head to host memory; decode speed is unchanged (228.5\,tok/s at
241{,}664 tokens with vision, 230.4 at 131{,}072 without).
These limits leave each engine's default free memory, 768\,MiB for azllm. With 2\,GB left free
for a desktop on the same card, azllm starts with at most 151{,}552 tokens with \texttt{q8/q4} KV
and 233{,}472 with \texttt{q4/q4}: near the \texttt{q8/q4} limit, what no longer fits is the
full-size prefill buffers, not the KV cache.

\begin{figure*}[t]
  \centering
  \begin{tikzpicture}
  \begin{axis}[
      xbar, width=0.66\textwidth, height=9.4cm, bar width=5pt,
      xmin=0, xmax=440, xlabel={largest context that starts (thousand tokens)},
      xlabel style={font=\small},
      symbolic y coords={oj,ojc,ni8m,ni8d,nr8m,nr8d,ne8m,ne8d,ne8n,a8,a84,a44},
      ytick=data, y dir=reverse,
      yticklabels={OpenJet q8 MTP 4, OpenJet q8 MTP 4 (512-token prefill chunks),
                   NInfer int8 MTP 3, NInfer int8 DFlash2 7, NInfer rk8v4 MTP 3,
                   NInfer rk8v4 DFlash2 7, NInfer E8 MTP 3, NInfer E8 DFlash2 7,
                   NInfer E8 no drafter, azllm q8 DFlash2 7, azllm q8/q4 DFlash2 7, azllm q4/q4 DFlash2 7},
      yticklabel style={font=\scriptsize}, xticklabel style={font=\scriptsize},
      legend style={font=\scriptsize, at={(0.5,1.02)}, anchor=south, draw=none,
                    legend columns=2, /tikz/every even column/.append style={column sep=8pt}},
      legend cell align=left,
      point meta=explicit symbolic,
      nodes near coords, nodes near coords style={font=\tiny},
      every axis plot/.append style={draw=none}, xmajorgrids, grid style={gray!20},
      extra x ticks={262.144}, extra x tick labels={native 262K},
      extra x tick style={grid=major, major grid style={cgood, thick, dashed},
                          tick label style={text=cgood, font=\scriptsize, yshift=-8pt}}]
    \addplot[fill=cnin] coordinates {(131.1,oj) [131K] (147.5,ojc) [147K] (172.0,ni8m) [172K]
        (114.7,ni8d) [115K] (229.4,nr8m) [229K] (151.6,nr8d) [152K] (335.9,ne8m) [336K]
        (225.3,ne8d) [225K] (409.6,ne8n) [410K]
        (184.3,a8) [184K] (254.0,a84) [254K] (262.1,a44) [262K+]};
    \addplot[fill=cnin2] coordinates {(0,oj) [] (0,ojc) [] (155.6,ni8m) [156K] (98.3,ni8d) [98K]
        (208.9,nr8m) [209K] (131.1,nr8d) [131K] (307.2,ne8m) [307K] (192.5,ne8d) [193K] (0,ne8n) []
        (155.6,a8) [156K] (241.7,a84) [242K] (262.1,a44) [262K+]};
    \legend{text only, with vision}
  \end{axis}
\end{tikzpicture}

  \caption{Largest context accepted at startup, binary-searched in steps of 4096 tokens. The
           baseline's own setting is 128{,}000 tokens. OpenJet has no vision; the no-drafter E8
           row is shown text only; azllm's search stopped at the native window. The dashed line
           marks the model's native window; going past it would need RoPE scaling, which we did
           not test.}
  \label{fig:capacity}
\end{figure*}

\paragraph{Does the context really fit?}
A configuration that starts may still fail later, when a long prompt allocates prefill buffers.
Table~\ref{tab:fill} fills each retained configuration to near its limit and queries it. Every
server answered every question and stayed up, and the vision-enabled ones still described an
image afterwards. The baseline slows to 55\,tok/s at 121K tokens; NInfer stays above
106\,tok/s at 252K, and azllm answers at 162\,tok/s with 252K tokens in context. azllm's prompt
cache also shows here: after reading 219K tokens in 147\,s, the next request, the same prompt
plus 15K tokens, was read in 15\,s.

\begin{table}[t]
  \centering
  \caption{Full-context runs: prompt filled with source code, twelve questions asking for the
           first 16 lines of a named function. ``Exact'': answers with all 16 lines in place;
           ``Similar'': mean line-sequence similarity; decode speed while answering.}
  \label{tab:fill}
  \small
  \setlength{\tabcolsep}{3pt}
  \begin{tabular}{@{}lrrrr@{}}
    \toprule
    Configuration & Prompt & Exact & Similar & tok/s \\
    \midrule
    Baseline, vision & 121{,}375 & 11/12$^a$ & 0.989 & 55.3 \\
    OpenJet q8, 147K & 142{,}807 & 12/12 & 1.000 & 125.1 \\
    NInfer fast, vision & 184{,}886 & 12/12 & 1.000 & 114.7 \\
    NInfer precise, vision & 196{,}772 & 12/12 & 1.000 & 109.9 \\
    NInfer precise, text & 221{,}095 & 11/12$^b$ & 0.995 & 106.3 \\
    NInfer long, vision & 252{,}392 & 12/12 & 1.000 & 108.4 \\
    azllm \texttt{q8/q4}, vision & 233{,}777 & 11/12$^a$ & 0.989 & 132.1 \\
    azllm \texttt{q4/q4}, vision & 252{,}392 & 12/12 & 1.000 & 161.8 \\
    \bottomrule
  \end{tabular}\\[2pt]
  \parbox{\columnwidth}{\scriptsize $^a$ The function contains the literal text
  \texttt{<|im\_start|>}, which llama.cpp and azllm read as a special token and drop from the
  output; the rest was exact. $^b$ The answer omitted the first line, \texttt{try \{}, and copied
  the other 15 exactly.}
\end{table}

\subsection{What a 4-bit KV cache costs on code}

The recall test of Table~\ref{tab:fill} does not separate the KV formats: at 98K tokens, int8,
rk8v4 and E8 all copy every requested line exactly, and the 4-bit formats of both NInfer and
azllm still do at 252K. Perplexity is more sensitive (Figure~\ref{fig:kvppl}). In NInfer, int8
is indistinguishable from bf16, rk8v4 costs 0.13\% (0.15\% at 128K context), and E8 with 4-bit
keys 0.82\% (0.79\%). In llama.cpp, the baseline's \texttt{turbo4} values cost 0.11\%. In
azllm, \texttt{q8/q4} costs 0.04\% and \texttt{q4/q4} 0.16\%, with a mean KL divergence from the
\texttt{q8} cache of 0.0008 and 0.0019 nats and top-1 agreement of 99.5\% and 99.1\%. azllm's
4-bit keys, rotated and quantized in blocks of 32 with one FP16 scale, are five times cheaper
than NInfer's $E_8$-lattice keys.

\begin{figure}[t]
  \centering
  \begin{tikzpicture}
  \begin{axis}[
      ybar, width=\columnwidth, height=5.4cm, bar width=11pt,
      ymin=-0.1, ymax=1.0, ytick={0,0.25,0.5,0.75,1},
      ylabel={$\Delta$ NLL per token (\%)}, ylabel style={font=\small},
      symbolic x coords={ni8,nr8,ne8,lt4,a84,a44},
      xtick={ni8,nr8,ne8,lt4,a84,a44}, xticklabels={NInfer\\int8, NInfer\\rk8v4, NInfer\\E8, llama.cpp\\q8/turbo4, azllm\\q8/q4, azllm\\q4/q4},
      xticklabel style={font=\scriptsize, align=center},
      yticklabel style={font=\scriptsize},
      enlarge x limits=0.1, clip=false,
      point meta=explicit symbolic, nodes near coords, nodes near coords style={font=\tiny},
      every axis plot/.append style={draw=none, bar shift=0pt}, ymajorgrids, grid style={gray!20}]
    \addplot[fill=cnin] coordinates {(ni8,-0.025) [$-$0.03] (nr8,0.126) [0.13] (ne8,0.823) [0.82]};
    \addplot[fill=cbase] coordinates {(lt4,0.113) [0.11]};
    \addplot[fill=cgood] coordinates {(a84,0.042) [0.04] (a44,0.163) [0.16]};
  \end{axis}
\end{tikzpicture}

  \caption{Cost of each KV format on code: relative change of the mean NLL per token against
           the most precise format of the same engine (NInfer: bf16; llama.cpp: q8\_0/q8\_0;
           azllm: q8), 64K windows with the last 32K positions scored. Lower is better.}
  \label{fig:kvppl}
\end{figure}
